\documentclass[conference]{IEEEtran}

% ============================================================
% CONTENT-SUPPORT PACKAGES
% ============================================================
\usepackage{amsmath,amssymb,amsthm}
\usepackage{mathtools}
\usepackage{graphicx}
\usepackage{algorithm}
\usepackage{algpseudocode}
\usepackage{tikz}
\usetikzlibrary{shapes,arrows,positioning}
\usepackage{shared/template_tgir_article}
\usepackage{booktabs}
\usepackage{longtable}
\usepackage{tabularx}
\usepackage{array}
\usepackage{enumitem}
\usepackage{xcolor}
\usepackage{listings}

\lstset{
    basicstyle=\ttfamily\small,
    breaklines=true,
    columns=fullflexible,
    frame=single
}

\newcommand{\product}{TGIR QuantLib Tools}
\newcommand{\code}[1]{\texttt{#1}}
\newcommand{\docdate}{2026-04-06}

% Visual design is controlled exclusively by template_tgir_article.sty and
% IEEEtran. Keep this file limited to content-support packages and commands.


% Theorem environments
\theoremstyle{definition}
\newtheorem{definition}{Definition}[section]
\newtheorem{proposition}{Proposition}[section]
\newtheorem{remark}{Remark}[section]

% Math commands
\newcommand{\E}{\mathbb{E}}
\newcommand{\Q}{\mathbb{Q}}
\newcommand{\R}{\mathbb{R}}
\newcommand{\N}{\mathbb{N}}
\newcommand{\ind}{\mathbf{1}}
\newcommand{\diff}{\mathrm{d}}
\newcommand{\LGD}{\mathrm{LGD}}
\newcommand{\PD}{\mathrm{PD}}
\newcommand{\EE}{\mathrm{EE}}
\newcommand{\CVA}{\mathrm{CVA}}
\newcommand{\PV}{\mathrm{PV}}
\newcommand{\NPV}{\mathrm{NPV}}

% Hyperref setup
\hypersetup{
    pdftitle={Exotic Derivatives Pricing: Callable Cross-Currency Swaps and Credit Contingent Swaps},
    pdfauthor={Tim Glauner, TG Investments and Research, LLC}
}

\title{Pricing of Complex Exotic Derivatives: Callable Cross-Currency Swaps and Credit Contingent Interest Rate Swaps}
\TGIRSetArticlePublicationDate{August 2026}
\author{\TGIRArticleAuthorBlock}

% ============================================================
% DOCUMENT
% ============================================================
\begin{document}
\maketitle

\begin{abstract}
This paper presents a comprehensive analysis of the pricing methodologies for two complex derivative structures: (1) a 40-year callable cross-currency interest rate swap (CCIRS) between EUR and USD with a zero-coupon fixed leg in USD versus €STR (ESTER) floating payments in EUR, featuring a single European-style exercise; and (2) a credit contingent interest rate swap (CCIRS) where an underlying 10-year vanilla fixed-versus-SOFR USD swap between a development bank and a tier-1 institution (JPM) is subject to third-party credit risk from a financial institution reference entity. We examine market-standard pricing approaches employed by tier-1 investment banks, discuss the theoretical frameworks including multi-factor Gaussian affine term structure models, hybrid interest rate-FX-credit models, and wrong-way risk considerations. We evaluate the capabilities and limitations of the QuantLib open-source library for pricing these instruments, identifying significant gaps in native functionality for production-grade implementations. Calibration methodologies for both products are discussed in detail.
\end{abstract}

\begin{IEEEkeywords}
Callable cross-currency swap, credit contingent swap, CVA, wrong-way risk,
Hull-White, G2++, LGM model, swaption calibration, QuantLib
\end{IEEEkeywords}

% ============================================================
\section{Introduction}
% ============================================================

The pricing of exotic interest rate and credit derivatives presents significant challenges due to the complex interactions between multiple risk factors. This paper analyzes two structurally distinct but equally challenging derivative products that represent the frontier of quantitative finance:

\begin{enumerate}
    \item \textbf{Callable Cross-Currency Swap (CCIRS)}: A 40-year deal exchanging USD zero-coupon fixed payments against €STR (ESTER) floating payments in EUR, with a single call option allowing one party to terminate the swap early.

    \item \textbf{Credit Contingent Interest Rate Swap (CC-IRS)}: A 10-year vanilla fixed-vs-SOFR USD swap between a development bank and JPMorgan, with embedded contingency on a third-party financial institution's credit status.
\end{enumerate}

Both products require hybrid modeling frameworks that capture the joint dynamics of multiple risk factors under a unified no-arbitrage pricing measure.

\subsection{Motivation and Market Context}

Cross-currency callable structures have gained prominence as corporations and sovereigns seek long-dated financing with embedded optionality. The 40-year tenor introduces substantial model risk, as small errors in volatility or mean-reversion calibration compound significantly over extended horizons.

Credit contingent swaps emerged from the XVA (credit/debit/funding valuation adjustments) revolution following the 2008 financial crisis. They represent explicit contractual treatment of what was previously embedded as bilateral credit risk, making the credit exposure a tradeable and hedgeable quantity.

\subsection{Paper Organization}

Section 2 provides detailed specifications for both deals. Section 3 develops the theoretical pricing framework for callable cross-currency swaps. Section 4 addresses credit contingent swap pricing and the associated wrong-way risk challenges. Section 5 presents calibration methodologies. Section 6 assesses QuantLib's capabilities for these products. Section 7 concludes with recommendations.

% ============================================================
\section{Deal Specifications}
% ============================================================

\subsection{Deal 1: EUR/USD Callable Cross-Currency Swap}

\begin{table}[ht]
\centering
\caption{Callable CCIRS Term Sheet}
\label{tab:ccirs_terms}
\begin{tabular}{ll}
\toprule
\textbf{Parameter} & \textbf{Specification} \\
\midrule
Trade Date & T \\
Effective Date & T + 2 business days \\
Maturity Date & T + 40 years \\
Call Option Style & European (single exercise) \\
Call Exercise Date & Typically year 5, 10, or 15 \\
\midrule
\multicolumn{2}{l}{\textbf{USD Leg (Fixed, Zero-Coupon)}} \\
Notional & USD 100,000,000 \\
Fixed Rate & $K_{\text{USD}}$ (to be solved) \\
Day Count & ACT/360 \\
Payment Frequency & Zero-coupon (single payment at maturity) \\
\midrule
\multicolumn{2}{l}{\textbf{EUR Leg (Floating, €STR)}} \\
Notional & EUR equivalent at initial FX rate \\
Floating Index & €STR (Euro Short-Term Rate) \\
Spread & $s_{\text{EUR}}$ basis points \\
Day Count & ACT/360 \\
Payment Frequency & Quarterly, compounded in arrears \\
\midrule
\multicolumn{2}{l}{\textbf{FX Exchanges}} \\
Initial Exchange & Yes (at trade inception) \\
Final Exchange & Yes (at maturity or exercise) \\
FX Rate Fixing & Spot rate at T \\
\midrule
\multicolumn{2}{l}{\textbf{Optionality}} \\
Option Holder & Payer of USD fixed (receiver of EUR float) \\
Exercise Type & European \\
Settlement & Physical (swap terminates) \\
\bottomrule
\end{tabular}
\end{table}

\subsubsection{Key Challenges}

\begin{enumerate}
    \item \textbf{Ultra-long tenor}: 40 years requires careful extrapolation of yield curves and volatility surfaces beyond liquid market observables.

    \item \textbf{Zero-coupon structure}: The entire USD fixed payment occurs at maturity, creating extreme duration and convexity exposure.

    \item \textbf{Cross-currency basis}: EUR/USD basis spreads must be modeled, particularly important for long-dated structures.

    \item \textbf{FX-rate correlation}: The callable feature depends on joint dynamics of EUR and USD rates plus EUR/USD FX rate.
\end{enumerate}

\subsection{Deal 2: Credit Contingent Interest Rate Swap}

\begin{table}[ht]
\centering
\caption{Credit Contingent IRS Term Sheet}
\label{tab:ccirs_credit_terms}
\begin{tabular}{ll}
\toprule
\textbf{Parameter} & \textbf{Specification} \\
\midrule
\multicolumn{2}{l}{\textbf{Counterparties}} \\
Party A & Development Bank (e.g., IBRD, EIB, KfW) \\
Party B & JPMorgan Chase \& Co. \\
Reference Entity & Third-party Financial Institution \\
\midrule
\multicolumn{2}{l}{\textbf{Underlying Swap}} \\
Trade Date & T \\
Effective Date & T + 2 business days \\
Maturity Date & T + 10 years \\
Notional & USD 250,000,000 \\
\midrule
\multicolumn{2}{l}{\textbf{Fixed Leg}} \\
Fixed Rate & $K$ (market rate at inception) \\
Day Count & 30/360 \\
Payment Frequency & Semi-annual \\
\midrule
\multicolumn{2}{l}{\textbf{Floating Leg}} \\
Floating Index & SOFR (Secured Overnight Financing Rate) \\
Spread & 0 bps \\
Day Count & ACT/360 \\
Payment Frequency & Quarterly, compounded in arrears \\
\midrule
\multicolumn{2}{l}{\textbf{Credit Contingency}} \\
Credit Event Trigger & Reference Entity Default \\
Credit Event Definition & ISDA 2014 Credit Derivatives Definitions \\
Settlement Type & Cash settlement \\
Contingent Payment & $\max(V_{\tau}, 0)$ or $V_{\tau}$ (depending on structure) \\
Recovery Rate & Market standard or auction-determined \\
\bottomrule
\end{tabular}
\end{table}

\subsubsection{Credit Contingent Structure Variants}

Two primary structures exist in practice:

\begin{definition}[Unilateral Credit Contingent Swap]
Upon default of the reference entity at time $\tau < T$, if the mark-to-market value of the underlying swap $V_{\tau}$ is positive from Party A's perspective, Party B must pay $V_{\tau}$ to Party A. The swap then terminates.
\end{definition}

\begin{definition}[Bilateral Credit Contingent Swap]
Upon default, the party for whom the swap has negative value must pay the absolute value $|V_{\tau}|$ to the other party. The swap terminates regardless of which party is in-the-money.
\end{definition}

% ============================================================
\section{Pricing Framework: Callable Cross-Currency Swaps}
% ============================================================

\subsection{General Pricing Principle}

The value of a callable CCIRS can be decomposed as:
\begin{equation}
V_{\text{Callable CCIRS}} = V_{\text{CCIRS}} - V_{\text{Swaption}}
\end{equation}
where $V_{\text{CCIRS}}$ is the value of the underlying non-callable cross-currency swap and $V_{\text{Swaption}}$ is the value of the embedded option to exercise early.

However, for cross-currency structures, this decomposition requires a \textit{hybrid model} capturing the joint dynamics of:
\begin{itemize}
    \item USD short rate $r^{\text{USD}}_t$
    \item EUR short rate $r^{\text{EUR}}_t$
    \item EUR/USD spot FX rate $X_t$
\end{itemize}

\subsection{Model Selection: Tier-1 Bank Practice}

\subsubsection{Industry Standards}

Tier-1 investment banks predominantly employ the following models for callable cross-currency products:

\begin{table}[ht]
\centering
\caption{Model Hierarchy for Callable CCIRS at Tier-1 Banks}
\label{tab:models}
\begin{tabularx}{\columnwidth}{>{\raggedright\arraybackslash}X>{\raggedright\arraybackslash}Xl}
\toprule
\textbf{Model} & \textbf{Use Case} & \textbf{Complexity} \\
\midrule
Hull-White 2-Factor (each ccy) & Standard callable CCIRS & Medium \\
Linear Gaussian Markov (LGM) & Production standard & Medium-High \\
G2++ per currency & Long-dated products & High \\
LIBOR Market Model (LMM) & Complex exotics & Very High \\
Multi-currency LMM & Exotic cross-currency & Extremely High \\
\bottomrule
\end{tabularx}
\end{table}

\subsubsection{The Linear Gaussian Markov (LGM) Model}

The LGM model, also known as the Hull-White one-factor model in its simplest form, has become the de facto standard for callable fixed income products at major banks due to its:

\begin{itemize}
    \item Analytical tractability for bond and European swaption pricing
    \item Efficient lattice (tree) or PDE implementation for Bermudan exercise
    \item Straightforward calibration to co-terminal swaption volatilities
\end{itemize}

The short rate dynamics under the risk-neutral measure $\Q$ are:
\begin{equation}
\diff r_t = [\theta(t) - a \cdot r_t] \diff t + \sigma(t) \diff W_t
\end{equation}
where:
\begin{itemize}
    \item $\theta(t)$ is calibrated to fit the initial term structure
    \item $a > 0$ is the mean-reversion speed
    \item $\sigma(t)$ is the time-dependent volatility (calibrated to swaptions)
\end{itemize}

\subsubsection{The G2++ Two-Factor Model}

For ultra-long-dated products (40 years), a two-factor model is often preferred to capture the term structure of volatility and correlation between short and long rates:

\begin{align}
r_t &= x_t + y_t + \phi(t) \\
\diff x_t &= -a \cdot x_t \diff t + \sigma_1 \diff W^1_t \\
\diff y_t &= -b \cdot y_t \diff t + \sigma_2 \diff W^2_t
\end{align}
where $\E[\diff W^1_t \cdot \diff W^2_t] = \rho \diff t$ and $\phi(t)$ ensures the model fits the initial yield curve.

\begin{proposition}[Two-Factor Advantage for Long Tenors]
The G2++ model produces more realistic swaption volatility term structures for long-dated options because:
\begin{enumerate}
    \item The correlation $\rho$ allows for humped or inverted volatility term structures
    \item Different mean-reversion speeds $a \neq b$ capture the decorrelation of short and long rates
    \item The model can fit both caplet and swaption volatilities simultaneously
\end{enumerate}
\end{proposition}

\subsection{Cross-Currency Extension}

For the EUR/USD CCIRS, we need a joint model. Under the domestic (USD) risk-neutral measure $\Q^{\text{USD}}$:

\begin{align}
\diff r^{\text{USD}}_t &= [\theta_{\text{USD}}(t) - a_{\text{USD}} r^{\text{USD}}_t] \diff t + \sigma_{\text{USD}} \diff W^{1}_t \\
\diff r^{\text{EUR}}_t &= [\theta_{\text{EUR}}(t) - a_{\text{EUR}} r^{\text{EUR}}_t] \diff t \nonumber\\
&\quad + \lambda_{\text{FX}} \sigma_{\text{EUR}} \sigma_X \rho_{r^{\text{EUR}},X}\diff t
 + \sigma_{\text{EUR}} \diff W^{2}_t \\
\diff X_t &= (r^{\text{USD}}_t - r^{\text{EUR}}_t) X_t \diff t + \sigma_X X_t \diff W^{3}_t
\end{align}

with correlation structure:
\begin{equation}
\text{Corr}(\diff W^i_t, \diff W^j_t) = \rho_{ij} \diff t, \quad i,j \in \{1,2,3\}
\end{equation}

\begin{remark}[Quanto Adjustment]
The drift adjustment term $\lambda_{\text{FX}} \sigma_{\text{EUR}} \sigma_X \rho_{r^{\text{EUR}},X}$ in the EUR rate dynamics arises from the change of measure from the EUR risk-neutral measure to the USD risk-neutral measure. This is the ``quanto drift'' and is essential for correct pricing of cross-currency derivatives.
\end{remark}

\subsection{Numerical Methods for Callable Structures}

\subsubsection{Lattice Methods (Trees)}

For single exercise (European) callable swaps, a recombining trinomial tree in the spirit of Hull-White can be used:

\begin{algorithm}
\caption{Trinomial Tree for Callable CCIRS}
\begin{algorithmic}[1]
\State Build two-dimensional tree for $(r^{\text{USD}}, r^{\text{EUR}})$ or use state variables
\State At terminal nodes: compute underlying swap value $V_{\text{swap}}$
\State Backward induction through tree:
\For{each time step from $T$ to $t_0$}
    \For{each node}
        \State $V_{\text{cont}} \gets e^{-r \Delta t} \E^{\Q}[V_{t+\Delta t} | \mathcal{F}_t]$
        \If{at exercise date}
            \State $V_t \gets \max(V_{\text{cont}}, 0)$ \Comment{Exercise if positive continuation}
        \Else
            \State $V_t \gets V_{\text{cont}}$
        \EndIf
    \EndFor
\EndFor
\State \Return $V_0$
\end{algorithmic}
\end{algorithm}

\subsubsection{Monte Carlo with Regression (LSM)}

For more complex structures or when the number of state variables is large, the Longstaff-Schwartz Least Squares Monte Carlo (LSM) algorithm is employed:

\begin{equation}
C(t_k, S_{t_k}) \approx \sum_{j=1}^{M} \beta_j \phi_j(S_{t_k})
\end{equation}

where $\phi_j$ are basis functions (typically polynomials) and $\beta_j$ are regression coefficients estimated from simulated paths.

% ============================================================
\section{Pricing Framework: Credit Contingent Swaps}
% ============================================================

\subsection{Conceptual Framework}

A credit contingent swap (CCS) combines interest rate swap pricing with credit derivative valuation. The key insight is that the product has embedded optionality triggered by a credit event.

\begin{definition}[Credit Contingent Swap Value]
Let $\tau$ denote the default time of the reference entity, $V_t$ the value of the underlying swap at time $t$, and $T$ the swap maturity. The credit contingent swap value from the protection buyer's perspective is:
\begin{multline}
V_{\text{CCS}} = \E^{\Q}\bigl[ \ind_{\{\tau \leq T\}} D(0,\tau) g(V_\tau) \\
{}+ \ind_{\{\tau > T\}} \text{(standard swap PV)} \bigr]
\end{multline}
where $g(V_\tau)$ defines the payoff structure upon default.
\end{definition}

\subsection{Connection to CVA}

The credit contingent swap is intimately related to Credit Valuation Adjustment (CVA). In fact, a unilateral CCS with payoff $g(V_\tau) = \max(V_\tau, 0)$ is equivalent to:

\begin{equation}
V_{\text{CCS}} = V_{\text{Swap}} - \CVA
\end{equation}

where:
\begin{equation}
\CVA = \LGD \int_0^T \EE^*(t) \cdot \diff \PD(0,t)
\end{equation}

with $\EE^*(t)$ being the discounted expected positive exposure and $\PD(0,t)$ the probability of default between time 0 and $t$.

\subsection{Wrong-Way Risk}

A critical consideration for financial institution reference entities is \textbf{wrong-way risk} (WWR): the tendency for exposure and default probability to be positively correlated.

\begin{definition}[Wrong-Way Risk]
Wrong-way risk exists when:
\begin{equation}
\text{Corr}(\text{Exposure at } \tau, \PD(\tau)) > 0
\end{equation}
This occurs when the reference entity is more likely to default precisely when the exposure is large.
\end{definition}

For a financial institution reference entity (such as another bank), WWR manifests because:
\begin{itemize}
    \item Financial stress periods (high rates, wide spreads) increase both swap exposure and default probability
    \item Systemic risk creates correlation between all financial institutions
    \item CDS spreads and interest rates tend to move together in stress scenarios
\end{itemize}

\subsection{Hybrid Credit-IR Models}

\subsubsection{Reduced-Form Credit Model}

The standard approach combines a Hull-White interest rate model with a Cox process for default intensity:

\begin{align}
\diff r_t &= [\theta(t) - a r_t] \diff t + \sigma_r \diff W^r_t \\
\diff \lambda_t &= \kappa(\bar{\lambda} - \lambda_t) \diff t + \sigma_\lambda \sqrt{\lambda_t} \diff W^\lambda_t \\
\E[\diff W^r_t \cdot \diff W^\lambda_t] &= \rho_{r,\lambda} \diff t
\end{align}

where $\lambda_t$ is the default intensity (hazard rate) and default occurs at the first jump of a Cox process with intensity $\lambda_t$.

\subsubsection{CIR++ Model for Default Intensity}

The shifted CIR model (CIR++) provides a flexible framework:

\begin{equation}
\lambda_t = y_t + \psi(t)
\end{equation}

where $y_t$ follows a standard CIR process and $\psi(t)$ is a deterministic shift calibrated to CDS market quotes.

\subsection{Pricing under Wrong-Way Risk}

\subsubsection{Independent Case (Baseline)}

Under independence, the CVA simplifies to:

\begin{equation}
\CVA_{\text{ind}} = \LGD \sum_{i=1}^{N} \EE^*(t_i) \cdot q(t_{i-1}, t_i)
\end{equation}

where $q(t_{i-1}, t_i)$ is the marginal default probability.

\subsubsection{Correlated Case (WWR)}

With wrong-way risk:

\begin{equation}
\CVA_{\text{WWR}} = \LGD \cdot \E^{\Q}\left[ D(0,\tau) \cdot V_\tau^+ \mid \tau \leq T \right] \cdot \PD(0,T)
\end{equation}

This requires joint simulation of interest rates and default intensity.

\subsubsection{Hull-White WWR Model}

A practical approach used by tier-1 banks:

\begin{equation}
\lambda_t = \lambda_0 \cdot \exp\left(\alpha \cdot \frac{V_t - V_0}{\sigma_V}\right)
\end{equation}

where $\alpha > 0$ captures the WWR effect: as exposure $V_t$ increases, so does the default intensity.

\subsection{Practical Implementation}

\begin{algorithm}
\caption{Monte Carlo CVA with Wrong-Way Risk}
\begin{algorithmic}[1]
\State Calibrate interest rate model to swaption vols
\State Calibrate credit model to CDS curve
\State Set correlation $\rho_{r,\lambda}$ (typically from historical analysis)
\For{$k = 1$ to $N_{\text{paths}}$}
    \State Simulate joint paths $(r_t^{(k)}, \lambda_t^{(k)})$ for $t \in [0,T]$
    \State Generate default time $\tau^{(k)}$ from $\lambda_t^{(k)}$
    \If{$\tau^{(k)} \leq T$}
        \State Compute swap value $V_{\tau^{(k)}}^{(k)}$
        \State $\text{Loss}^{(k)} \gets \LGD \cdot D(0,\tau^{(k)}) \cdot \max(V_{\tau^{(k)}}^{(k)}, 0)$
    \Else
        \State $\text{Loss}^{(k)} \gets 0$
    \EndIf
\EndFor
\State \Return $\CVA \gets \frac{1}{N_{\text{paths}}} \sum_{k=1}^{N_{\text{paths}}} \text{Loss}^{(k)}$
\end{algorithmic}
\end{algorithm}

% ============================================================
\section{Model Calibration}
% ============================================================

\subsection{Calibration for Callable CCIRS}

\subsubsection{Yield Curve Construction}

For USD and EUR curves in 2026:
\begin{itemize}
    \item \textbf{USD}: SOFR OIS curve from deposits, futures, and swaps
    \item \textbf{EUR}: €STR OIS curve from deposits and swaps
    \item \textbf{Cross-currency basis}: EUR/USD basis swap curve
\end{itemize}

\subsubsection{Swaption Volatility Calibration}

The Hull-White/LGM model is typically calibrated to co-terminal European swaptions:

\begin{table}[ht]
\centering
\caption{Example Calibration Instruments for 40Y Callable Swap with 10Y Call}
\label{tab:calib_inst}
\begin{tabularx}{\columnwidth}{>{\raggedright\arraybackslash}X>{\raggedright\arraybackslash}Xl}
\toprule
\textbf{Instrument} & \textbf{Expiry} & \textbf{Tenor} \\
\midrule
EUR Swaption & 10Y & 30Y \\
EUR Swaption & 10Y & 25Y \\
EUR Swaption & 10Y & 20Y \\
EUR Swaption & 10Y & 15Y \\
EUR Swaption & 10Y & 10Y \\
\midrule
USD Swaption & 10Y & 30Y \\
USD Swaption & 10Y & 25Y \\
\multicolumn{3}{c}{...} \\
\bottomrule
\end{tabularx}
\end{table}

The objective function minimizes:
\begin{equation}
\min_{\{a, \sigma(t)\}} \sum_{i} w_i \left( \sigma^{\text{model}}_i(a, \sigma(\cdot)) - \sigma^{\text{market}}_i \right)^2
\end{equation}

\subsubsection{Mean Reversion Calibration}

For long-dated products, mean reversion $a$ is crucial:
\begin{itemize}
    \item Too high: excessive dampening of long-rate volatility
    \item Too low: unrealistic long-dated vol explosion
    \item Typical range: $a \in [0.01, 0.05]$ (1-5\% annual mean reversion)
\end{itemize}

\subsection{Calibration for Credit Contingent Swap}

\subsubsection{CDS Curve Bootstrapping}

The hazard rate term structure is calibrated from CDS quotes:

\begin{equation}
\text{CDS Spread}_T = \frac{(1-R) \int_0^T D(0,t) \lambda(t) S(t) \diff t}{\sum_{i} \delta_i D(0,t_i) S(t_i)}
\end{equation}

where $S(t) = \exp(-\int_0^t \lambda(s) \diff s)$ is the survival probability.

\subsubsection{Correlation Estimation}

The interest rate-credit correlation $\rho_{r,\lambda}$ can be estimated from:
\begin{itemize}
    \item Historical correlation between CDS spreads and swap rates
    \item Implied from swaption-CDS basis trades
    \item Expert judgment based on market conditions
\end{itemize}

For financial institution reference entities, typical values:
\begin{equation}
\rho_{r,\lambda} \in [0.2, 0.5]
\end{equation}

Higher values in stress periods, lower in benign markets.

% ============================================================
\section{QuantLib Assessment}
% ============================================================

\subsection{Overview of QuantLib Capabilities}

QuantLib is a comprehensive open-source library for quantitative finance written in C++ with Python bindings. We assess its suitability for pricing the two deal types.

\subsection{Deal 1: Callable Cross-Currency Swap}

\subsubsection{Available Components}

\begin{table*}[!t]
\centering
\caption{QuantLib Components for Callable CCIRS}
\label{tab:ql_ccirs}
\begin{tabular}{lll}
\toprule
\textbf{Component} & \textbf{QuantLib Class} & \textbf{Status} \\
\midrule
Hull-White 1F Model & \texttt{HullWhite} & \checkmark Available \\
G2++ Model & \texttt{G2} & \checkmark Available \\
Bermudan Swaption & \texttt{Swaption} & \checkmark Available \\
Tree Engine & \texttt{TreeSwaptionEngine} & \checkmark Available \\
FD Engine & \texttt{FdHullWhiteSwaptionEngine} & \checkmark Available \\
Cross-Currency Swap & \texttt{CrossCurrencyBasisSwap} & Partial \\
Callable CCIRS & -- & \textbf{Not Available} \\
Multi-currency Model & -- & \textbf{Not Available} \\
\bottomrule
\end{tabular}
\end{table*}

\subsubsection{Critical Gaps}

\begin{enumerate}
    \item \textbf{No native multi-currency model}: QuantLib lacks a built-in framework for correlated multi-currency short-rate dynamics under a single measure.

    \item \textbf{No callable cross-currency instrument}: The \texttt{CallableFixedRateBond} exists, but no equivalent for CCIRS.

    \item \textbf{Quanto drift handling}: Manual implementation required for measure changes.

    \item \textbf{FX-rate correlation}: Must be handled externally.
\end{enumerate}

\subsubsection{Implementation Path}

A production implementation would require:

\begin{lstlisting}[language=Python, caption=Conceptual QuantLib Extension]
# Pseudo-code for QuantLib extension
class MultiCurrencyHullWhite:
    def __init__(self, curves, vols, correlations):
        self.hw_usd = ql.HullWhite(curves['USD'])
        self.hw_eur = ql.HullWhite(curves['EUR'])
        self.fx_vol = vols['EURUSD']
        self.correlations = correlations

    def simulate(self, dates, paths):
        # Generate correlated Brownian motions
        # Apply quanto drift adjustment
        # Return joint paths for (r_USD, r_EUR, X)
        pass

    def price_callable_ccirs(self, swap, exercise_dates):
        # LSM or tree-based pricing
        pass
\end{lstlisting}

\subsubsection{Effort Estimate}

\begin{itemize}
    \item Development: 3-6 months for experienced quant developer
    \item Testing and validation: 2-3 months
    \item Production hardening: 2-3 months
\end{itemize}

\subsection{Deal 2: Credit Contingent Swap}

\subsubsection{Available Components}

\begin{table*}[!t]
\centering
\caption{QuantLib Components for Credit Contingent Swap}
\label{tab:ql_ccs}
\begin{tabular}{lll}
\toprule
\textbf{Component} & \textbf{QuantLib Class} & \textbf{Status} \\
\midrule
Vanilla IRS & \texttt{VanillaSwap} & \checkmark Available \\
CDS Pricing & \texttt{CreditDefaultSwap} & \checkmark Available \\
Hazard Rate Curve & \texttt{DefaultProbabilityTermStructure} & \checkmark Available \\
Monte Carlo Engine & \texttt{MCVanillaEngine} & \checkmark Available \\
CVA Calculation & -- & \textbf{Not Native} \\
Wrong-Way Risk & -- & \textbf{Not Available} \\
Credit-IR Correlation & -- & \textbf{Not Available} \\
Credit Contingent Swap & -- & \textbf{Not Available} \\
\bottomrule
\end{tabular}
\end{table*}

\subsubsection{Critical Gaps}

\begin{enumerate}
    \item \textbf{No CVA framework}: QuantLib does not provide native CVA/DVA/FVA calculations.

    \item \textbf{No hybrid credit-IR model}: Must implement joint dynamics externally.

    \item \textbf{No exposure simulation}: EPE/ENE calculation requires custom code.

    \item \textbf{No wrong-way risk}: Correlation between rates and credit must be implemented.
\end{enumerate}

\subsubsection{Possible Workaround}

Using QuantLib building blocks:

\begin{lstlisting}[language=Python, caption=CVA Approximation with QuantLib]
import QuantLib as ql
import numpy as np

def approximate_cva(swap, cds_curve, lgd, dates):
    """
    Simplified CVA calculation using QuantLib components
    WARNING: Does not capture wrong-way risk
    """
    cva = 0.0
    for i in range(1, len(dates)):
        # Expected exposure (simplified - should use simulation)
        npv = swap.NPV()
        ee = max(npv, 0)  # Crude approximation

        # Default probability
        surv_prev = cds_curve.survivalProbability(dates[i-1])
        surv_curr = cds_curve.survivalProbability(dates[i])
        pd = surv_prev - surv_curr

        # Discount factor
        df = swap.discountCurve().discount(dates[i])

        cva += lgd * df * ee * pd

    return cva
\end{lstlisting}

\subsection{Overall Assessment}

\begin{table}[ht]
\centering
\caption{QuantLib Suitability Summary}
\label{tab:ql_summary}
\begin{tabularx}{\columnwidth}{l>{\centering\arraybackslash}X>{\centering\arraybackslash}X>{\centering\arraybackslash}X}
\toprule
\textbf{Criterion} & \textbf{Callable CCIRS} & \textbf{Credit Contingent} & \textbf{Tier-1 Standard} \\
\midrule
Native Support & No & No & Yes \\
Building Blocks & Partial & Partial & Yes \\
Development Effort & High & Medium-High & N/A \\
Production Ready & No & No & Yes \\
Research/POC Use & Possible & Possible & N/A \\
\bottomrule
\end{tabularx}
\end{table}

\begin{remark}[QuantLib Limitations for Production]
While QuantLib provides excellent building blocks for interest rate modeling, it is not designed for XVA calculations or complex multi-asset hybrid models. Tier-1 banks typically use proprietary systems (e.g., Bloomberg MARS, Numerix, TriOptima) or in-house platforms for these products.
\end{remark}

\subsection{Alternative Open-Source Options}

\begin{itemize}
    \item \textbf{ORE (Open Source Risk Engine)}: Extension of QuantLib with XVA capabilities
    \item \textbf{Strata (OpenGamma)}: Java-based library with some XVA support
    \item \textbf{TF Quant Finance}: Google's TensorFlow-based library (GPU acceleration)
\end{itemize}

% ============================================================
\section{Conclusion}
% ============================================================

\subsection{Summary of Findings}

\begin{enumerate}
    \item \textbf{Callable CCIRS Pricing}: Requires multi-factor Gaussian models (G2++ or multi-currency LGM) with careful calibration to long-dated swaption volatilities. Tier-1 banks use hybrid models with FX-rate correlation. QuantLib lacks native support but provides useful building blocks.

    \item \textbf{Credit Contingent Swap Pricing}: Requires hybrid credit-IR models with explicit wrong-way risk treatment. CVA frameworks are essential. QuantLib is insufficient for production use; ORE or proprietary systems are necessary.

    \item \textbf{Model Calibration}: Both products require sophisticated calibration routines. For long-dated structures, mean reversion and volatility term structure are critical. For credit products, the interest rate-credit correlation drives wrong-way risk magnitude.
\end{enumerate}

\subsection{Recommendations}

\begin{enumerate}
    \item For \textbf{research and prototyping}, QuantLib can be extended with custom multi-currency and CVA modules.

    \item For \textbf{production pricing}, consider:
    \begin{itemize}
        \item ORE (Open Source Risk Engine) for XVA
        \item Commercial vendors (Numerix, FINCAD, Bloomberg)
        \item In-house development with significant investment
    \end{itemize}

    \item For \textbf{model validation}, compare multiple approaches:
    \begin{itemize}
        \item Analytical approximations where available
        \item Monte Carlo with variance reduction
        \item PDE/Lattice methods for cross-validation
    \end{itemize}
\end{enumerate}

\subsection{Future Research Directions}

\begin{itemize}
    \item Machine learning calibration for multi-factor models
    \item Neural network-based American/Bermudan exercise boundary approximation
    \item Deep hedging approaches for exotic cross-currency products
    \item Real-time CVA calculation with GPU acceleration
\end{itemize}

% ============================================================
% REFERENCES
% ============================================================
\section*{References}
\addcontentsline{toc}{section}{References}

\begin{thebibliography}{99}

\bibitem{andersen2010}
Andersen, L. and Piterbarg, V. (2010). \textit{Interest Rate Modeling}. Atlantic Financial Press.

\bibitem{brigo2006}
Brigo, D. and Mercurio, F. (2006). \textit{Interest Rate Models - Theory and Practice}. Springer, 2nd edition.

\bibitem{gregory2015}
Gregory, J. (2015). \textit{The xVA Challenge: Counterparty Credit Risk, Funding, Collateral and Capital}. Wiley.

\bibitem{hull2017}
Hull, J. and White, A. (2017). ``Optimal Delta Hedging for Options,'' \textit{Journal of Banking and Finance}, 82, 180-190.

\bibitem{hull1990}
Hull, J. and White, A. (1990). ``Pricing Interest-Rate Derivative Securities,'' \textit{Review of Financial Studies}, 3(4), 573-592.

\bibitem{jamshidian1989}
Jamshidian, F. (1989). ``An Exact Bond Option Formula,'' \textit{Journal of Finance}, 44(1), 205-209.

\bibitem{longstaff2001}
Longstaff, F. and Schwartz, E. (2001). ``Valuing American Options by Simulation: A Simple Least-Squares Approach,'' \textit{Review of Financial Studies}, 14(1), 113-147.

\bibitem{pykhtin2007}
Pykhtin, M. and Zhu, S. (2007). ``A Guide to Modeling Counterparty Credit Risk,'' \textit{GARP Risk Review}, July/August.

\bibitem{rebonato2004}
Rebonato, R. (2004). \textit{Volatility and Correlation}. Wiley, 2nd edition.

\bibitem{shreve2004}
Shreve, S. (2004). \textit{Stochastic Calculus for Finance II: Continuous-Time Models}. Springer.

\bibitem{turfus2020}
Turfus, C. (2020). ``Caplet Pricing with Backward-Looking Rates,'' SSRN Working Paper.

\end{thebibliography}

% ============================================================
% APPENDIX
% ============================================================
\begin{appendices}

\section{Mathematical Derivations}

\subsection{Hull-White Bond Price Formula}

Under the Hull-White model, the time-$t$ price of a zero-coupon bond maturing at $T$ is:
\begin{equation}
P(t,T) = A(t,T) \exp(-B(t,T) r_t)
\end{equation}
where:
\begin{align}
B(t,T) &= \frac{1 - e^{-a(T-t)}}{a}, \\
A(t,T) &= \frac{P(0,T)}{P(0,t)} \exp\biggl(
 B(t,T) \frac{\partial \ln P(0,t)}{\partial t} \nonumber\\
&\qquad - \frac{\sigma^2}{4a}(1-e^{-2at})B(t,T)^2 \biggr).
\end{align}

\subsection{G2++ Swaption Pricing}

The value of a European payer swaption in the G2++ model can be approximated using Jamshidian's decomposition, treating the swap as a portfolio of zero-coupon bonds and applying the bond option formula.

\subsection{CVA Formula Derivation}

Starting from the definition:
\begin{equation}
\CVA = \E^{\Q}\left[ \LGD \cdot D(0,\tau) \cdot V_\tau^+ \cdot \ind_{\tau \leq T} \right]
\end{equation}

Under independence of exposure and default:
\begin{align}
\CVA &= \LGD \int_0^T \E^{\Q}\left[ D(0,t) V_t^+ \right] f_\tau(t) \diff t \\
&= \LGD \int_0^T \EE^*(t) \cdot \lambda(t) S(t) \diff t \\
&= \LGD \int_0^T \EE^*(t) \diff(1 - S(t)) \\
&= \LGD \int_0^T \EE^*(t) \diff \PD(0,t)
\end{align}

\section{QuantLib Code Examples}

\subsection{Hull-White Calibration}

\begin{lstlisting}[language=Python, caption=Hull-White Calibration in QuantLib]
import QuantLib as ql

# Setup
today = ql.Date(14, 8, 2026)
ql.Settings.instance().evaluationDate = today

# Yield curve (SOFR OIS)
sofr_curve = ql.FlatForward(today, 0.04, ql.Actual365Fixed())
sofr_handle = ql.YieldTermStructureHandle(sofr_curve)

# Hull-White model
hw_model = ql.HullWhite(sofr_handle)

# Calibration helpers (swaptions)
swaption_helpers = []
for expiry, tenor, vol in calibration_data:
    helper = ql.SwaptionHelper(
        ql.Period(expiry),
        ql.Period(tenor),
        ql.QuoteHandle(ql.SimpleQuote(vol)),
        sofr_index,
        ql.Period(1, ql.Years),
        ql.Thirty360(),
        ql.Actual360(),
        sofr_handle
    )
    helper.setPricingEngine(
        ql.JamshidianSwaptionEngine(hw_model)
    )
    swaption_helpers.append(helper)

# Calibrate
optimization_method = ql.LevenbergMarquardt()
end_criteria = ql.EndCriteria(1000, 100, 1e-8, 1e-8, 1e-8)
hw_model.calibrate(swaption_helpers, optimization_method, end_criteria)

print(f"Mean reversion: {hw_model.params()[0]:.4f}")
print(f"Volatility: {hw_model.params()[1]:.4f}")
\end{lstlisting}

\subsection{CDS Curve Bootstrap}

\begin{lstlisting}[language=Python, caption=CDS Curve Bootstrap in QuantLib]
import QuantLib as ql

# CDS quotes (spreads in bps)
cds_quotes = [
    (ql.Period(1, ql.Years), 50),
    (ql.Period(3, ql.Years), 75),
    (ql.Period(5, ql.Years), 100),
    (ql.Period(7, ql.Years), 120),
    (ql.Period(10, ql.Years), 140),
]

# Build helpers
helpers = []
for tenor, spread in cds_quotes:
    helper = ql.SpreadCdsHelper(
        ql.QuoteHandle(ql.SimpleQuote(spread / 10000)),
        tenor,
        0,  # settlement days
        ql.TARGET(),
        ql.Quarterly,
        ql.Following,
        ql.DateGeneration.TwentiethIMM,
        ql.Actual360(),
        0.4,  # recovery rate
        sofr_handle
    )
    helpers.append(helper)

# Bootstrap hazard rate curve
hazard_curve = ql.PiecewiseFlatHazardRate(
    today,
    helpers,
    ql.Actual365Fixed()
)

# Survival probabilities
for years in [1, 3, 5, 7, 10]:
    date = today + ql.Period(years, ql.Years)
    surv = hazard_curve.survivalProbability(date)
    print(f"{years}Y survival prob: {surv:.4%}")
\end{lstlisting}

\end{appendices}

\end{document}
